% Generated by experiments/scripts/export_latex.py (export_latex_c12.py) -- do not edit by hand.
% Source: experiments/data/2-processed-tables-and-evaluations/revised-contract-12m/discovery/discovery_tobs_both.json
\begin{table}[htbp]
\centering
\caption[발견 순서 검증: 판정 대비의 민감도.]{$T_{\mathrm{obs}}$의 주 방식에서 본 발견 순서 검증이에요. 두 판정 대비 $D_1$(C-PR$'$ 빼기 BFS)과 $D_2$(C-PR$'$ 빼기 $\mathrm{IB}'_U$)의 민감도를 보여 줘요. 칸마다 값 아래에, 뜨거운 계약들을 짝지어 다시 뽑은 부트스트랩의 95\% 백분위 구간(다시 뽑기 400번, 시드 42)을 적었어요. 행마다 설정을 하나씩만 바꿔요. 바꾸는 것은 면적을 재는 우주의 몫 $x$, 뜨거운 집합의 몫 $q$, 허브(합집합 입차수가 가장 큰 계약들, 뜨거운 집합에서는 빼지만 그래프에는 남김)를 뺀 뜨거운 집합, 1달러 이상인 사용자 계정의 수로 잰 중요도, 뜨거운 계약 대신 출발 계약을 다시 뽑는 부트스트랩, 그리고 첫 출발 계약들에서 순서 규칙을 네 배 자주 다시 계산하는 일정이에요. 우주에서 $I_A$가 0보다 큰 계약은 11{,}662개뿐이에요. 그래서 $I_A$의 뜨거운 집합은 두 몫 $q$에서 똑같아요. 다른 방식과 다른 동결은 표~\ref{tab:discovery-decision}에 있어요. 이 행들에는 판정이 없어요. 발견 순서 검증은 등록 분석의 결과를 안 뒤에 정했고, 계획과 등록 문서의 어떤 규칙도 정하지 않아요.}
\label{tab:discovery-sensitivity}
\fitwidth{%
\begin{tabular}{lcccc}
\toprule
 & \multicolumn{2}{c}{$H_q(I_T)$에서} & \multicolumn{2}{c}{$H_q(I_A)$에서} \\
\cmidrule(lr){2-3} \cmidrule(lr){4-5}
설정 & $D_1$ & $D_2$ & $D_1$ & $D_2$ \\
\midrule
판정 ($x=10\%$, $q=1\%$) & \begin{tabular}[c]{@{}c@{}}+0.0520\\[-1pt]{\footnotesize [+0.0485, +0.0558]}\end{tabular} & \begin{tabular}[c]{@{}c@{}}-0.0051\\[-1pt]{\footnotesize [-0.0087, -0.0010]}\end{tabular} & \begin{tabular}[c]{@{}c@{}}+0.0061\\[-1pt]{\footnotesize [+0.0014, +0.0111]}\end{tabular} & \begin{tabular}[c]{@{}c@{}}-0.0571\\[-1pt]{\footnotesize [-0.0621, -0.0519]}\end{tabular} \\
\midrule
$x=5\%$ & \begin{tabular}[c]{@{}c@{}}+0.0673\\[-1pt]{\footnotesize [+0.0635, +0.0714]}\end{tabular} & \begin{tabular}[c]{@{}c@{}}+0.0026\\[-1pt]{\footnotesize [-0.0010, +0.0064]}\end{tabular} & \begin{tabular}[c]{@{}c@{}}+0.0472\\[-1pt]{\footnotesize [+0.0429, +0.0521]}\end{tabular} & \begin{tabular}[c]{@{}c@{}}-0.0132\\[-1pt]{\footnotesize [-0.0178, -0.0084]}\end{tabular} \\
$x=20\%$ & \begin{tabular}[c]{@{}c@{}}+0.0199\\[-1pt]{\footnotesize [+0.0169, +0.0230]}\end{tabular} & \begin{tabular}[c]{@{}c@{}}-0.0063\\[-1pt]{\footnotesize [-0.0091, -0.0034]}\end{tabular} & \begin{tabular}[c]{@{}c@{}}-0.0518\\[-1pt]{\footnotesize [-0.0575, -0.0460]}\end{tabular} & \begin{tabular}[c]{@{}c@{}}-0.0966\\[-1pt]{\footnotesize [-0.1023, -0.0905]}\end{tabular} \\
$q=5\%$ & \begin{tabular}[c]{@{}c@{}}+0.0132\\[-1pt]{\footnotesize [+0.0121, +0.0144]}\end{tabular} & \begin{tabular}[c]{@{}c@{}}-0.0067\\[-1pt]{\footnotesize [-0.0078, -0.0057]}\end{tabular} & \begin{tabular}[c]{@{}c@{}}+0.0061\\[-1pt]{\footnotesize [+0.0014, +0.0111]}\end{tabular} & \begin{tabular}[c]{@{}c@{}}-0.0571\\[-1pt]{\footnotesize [-0.0621, -0.0519]}\end{tabular} \\
허브 1{,}170개를 뺌 & \begin{tabular}[c]{@{}c@{}}+0.0306\\[-1pt]{\footnotesize [+0.0274, +0.0341]}\end{tabular} & \begin{tabular}[c]{@{}c@{}}-0.0126\\[-1pt]{\footnotesize [-0.0159, -0.0088]}\end{tabular} & \begin{tabular}[c]{@{}c@{}}-0.0019\\[-1pt]{\footnotesize [-0.0065, +0.0026]}\end{tabular} & \begin{tabular}[c]{@{}c@{}}-0.0607\\[-1pt]{\footnotesize [-0.0654, -0.0553]}\end{tabular} \\
수로 잰 중요도 ($N_T$, $N_A$) & \begin{tabular}[c]{@{}c@{}}+0.0188\\[-1pt]{\footnotesize [+0.0163, +0.0211]}\end{tabular} & \begin{tabular}[c]{@{}c@{}}-0.0282\\[-1pt]{\footnotesize [-0.0308, -0.0254]}\end{tabular} & \begin{tabular}[c]{@{}c@{}}+0.0191\\[-1pt]{\footnotesize [+0.0129, +0.0262]}\end{tabular} & \begin{tabular}[c]{@{}c@{}}-0.0473\\[-1pt]{\footnotesize [-0.0537, -0.0408]}\end{tabular} \\
출발 계약을 다시 뽑는 부트스트랩 & \begin{tabular}[c]{@{}c@{}}+0.0520\\[-1pt]{\footnotesize [+0.0320, +0.0717]}\end{tabular} & \begin{tabular}[c]{@{}c@{}}-0.0051\\[-1pt]{\footnotesize [-0.0249, +0.0179]}\end{tabular} & \begin{tabular}[c]{@{}c@{}}+0.0061\\[-1pt]{\footnotesize [-0.0315, +0.0398]}\end{tabular} & \begin{tabular}[c]{@{}c@{}}-0.0571\\[-1pt]{\footnotesize [-0.0873, -0.0221]}\end{tabular} \\
$\rho/4$ (1/200만큼 늘 때마다 다시 계산, 출발 계약 5개) & \begin{tabular}[c]{@{}c@{}}+0.0381\\[-1pt]{\footnotesize [+0.0347, +0.0418]}\end{tabular} & \begin{tabular}[c]{@{}c@{}}-0.0129\\[-1pt]{\footnotesize [-0.0167, -0.0087]}\end{tabular} & \begin{tabular}[c]{@{}c@{}}-0.0134\\[-1pt]{\footnotesize [-0.0181, -0.0082]}\end{tabular} & \begin{tabular}[c]{@{}c@{}}-0.0713\\[-1pt]{\footnotesize [-0.0765, -0.0658]}\end{tabular} \\
\bottomrule
\end{tabular}
}
\end{table}
